% Build from paper/. Only real downloaded tuning aggregates enter this file.
\documentclass[11pt]{article}
\usepackage[a4paper,margin=2.4cm]{geometry}
\usepackage[T1]{fontenc}
\usepackage[utf8]{inputenc}
\usepackage{tgpagella,booktabs,graphicx,xcolor,hyperref,tikz}
% Shared manuscript and TikZ visual language for the OPERA paper.
% Keep colours and Helvetica-like figure typography synchronized with
% opera/visualization/style.py (Arial/Liberation Sans on render hosts).

\definecolor{OperaIndigo}{HTML}{2D3A8C}
\definecolor{OperaNavy}{HTML}{1A237E}
\definecolor{OperaPeriwinkle}{HTML}{5C7AEA}
\definecolor{OperaSlate}{HTML}{737D8C}
\definecolor{OperaTeal}{HTML}{408579}
\definecolor{OperaGreen}{HTML}{1B7A4A}
\definecolor{OperaSienna}{HTML}{B5451B}
\definecolor{OperaTerracotta}{HTML}{C16A58}
\definecolor{OperaOchre}{HTML}{B18B4D}
\definecolor{OperaLavender}{HTML}{8974AC}
\definecolor{OperaInk}{HTML}{17212B}
\definecolor{OperaInkSecondary}{HTML}{485463}
\definecolor{OperaInkMuted}{HTML}{657180}
\definecolor{OperaBorder}{HTML}{CBD2DA}
\definecolor{OperaGrid}{HTML}{E8ECF1}
\definecolor{OperaPaper}{HTML}{F7F8FA}

\tikzset{
  opera/.style={font=\sffamily\small, text=OperaInk},
  panel/.style={draw=none, fill=white, inner sep=0pt},
  softpanel/.style={
    rounded corners=0.8mm,
    draw=OperaGrid,
    line width=0.45pt,
    fill=OperaPaper,
    inner sep=2.4mm,
  },
  stage/.style={
    rounded corners=0.7mm,
    draw=OperaBorder,
    line width=0.6pt,
    fill=white,
    align=center,
    minimum height=8mm,
    inner xsep=3mm,
    inner ysep=1.8mm,
  },
  stageblue/.style={stage, draw=OperaIndigo!55, fill=OperaIndigo!7},
  stageteal/.style={stage, draw=OperaTeal!65, fill=OperaTeal!8},
  stagewarm/.style={stage, draw=OperaSienna!60, fill=OperaSienna!7},
  latent/.style={
    circle,
    draw=OperaIndigo!65,
    fill=OperaIndigo!9,
    line width=0.6pt,
    minimum size=6mm,
    inner sep=0pt,
    font=\sffamily\scriptsize\bfseries,
  },
  token/.style={
    rounded corners=0.7mm,
    draw=OperaBorder,
    fill=white,
    minimum width=6.2mm,
    minimum height=4.8mm,
    inner sep=0.5mm,
    font=\sffamily\tiny,
    align=center,
  },
  outcome/.style={
    rounded corners=1mm,
    draw=OperaTerracotta!70,
    fill=OperaTerracotta!9,
    minimum height=5.8mm,
    inner xsep=1.8mm,
    font=\sffamily\scriptsize,
    align=center,
  },
  flow/.style={
    -{Latex[length=2mm,width=1.2mm]},
    line width=0.75pt,
    draw=OperaInkSecondary,
  },
  flowblue/.style={flow, draw=OperaIndigo},
  flowteal/.style={flow, draw=OperaTeal},
  flowwarm/.style={flow, draw=OperaSienna},
  relation/.style={
    -{Latex[length=1.8mm,width=1.1mm]},
    line width=0.65pt,
    draw=OperaLavender,
    densely dashed,
  },
  captionnote/.style={
    font=\sffamily\fontsize{6.6}{7.5}\selectfont,
    text=OperaInkMuted,
    align=center,
  },
  labelsmall/.style={
    font=\sffamily\fontsize{7}{8}\selectfont,
    text=OperaInkSecondary,
    align=center,
  },
  heading/.style={
    font=\sffamily\fontsize{9}{10}\selectfont\bfseries,
    text=OperaInk,
    align=left,
  },
  panelletter/.style={font=\sffamily\bfseries\normalsize, text=OperaInk},
  divider/.style={draw=OperaBorder, line width=0.45pt},
  eyebrow/.style={font=\sffamily\fontsize{6.5}{7.5}\selectfont\bfseries,
    text=OperaIndigo, align=left},
}

\newcommand{\opera}{\textsc{opera}}
\newcommand{\result}[1]{\csname opera@res@#1\endcsname}
\input{../scratch/paper_reading_20261028/deep_rehearsal_review_20261028/manuscript_numbers/results_macros.tex}
\setlength{\parindent}{0pt}
\setlength{\parskip}{6pt}
\newcommand{\auditfigure}[2]{\clearpage\section*{#2}
  \IfFileExists{../scratch/paper_reading_20261028/graphical_audit_20261028/ready_for_audit/#1.pdf}{%
  \includegraphics[width=\textwidth,height=.82\textheight,keepaspectratio]{../scratch/paper_reading_20261028/graphical_audit_20261028/ready_for_audit/#1.pdf}}{}
}
\AddToHook{shipout/background}{\put(0.5\paperwidth,-0.9cm){\makebox[0pt]{\sffamily\scriptsize\color{teal}
  TUNING SNAPSHOT --- development evidence, not held-out test findings}}}
\title{OPERA tuning results: audit snapshot}
\author{For manuscript review}
\date{28 October 2026}
\begin{document}
\maketitle
\textbf{Scope.} This is a factual results companion to the full manuscript.
It uses the inspected server bundle dated 28 October 2026. No national-model
results or test findings are included. Missing analyses are listed explicitly.
Synthetic layout checks are excluded. Figure designs remain subject to the
separate graphical change log.

\section*{Available local results}
The principal equal-cell summaries contain 3,004 supported
disease--outcome--horizon comparisons. One-year comparisons use a common
roster of 600 cells. Pooling horizons within each disease--outcome cell gives
675 cells. These denominators answer different questions and are not
interchangeable. Supported comparisons require at least five cases and five
controls. The current implementation applies that roster to Brier and IPA
as well as AUC.

\begin{center}\small
\begin{tabular}{lr}\toprule
Comparison & AUC difference, 95\% CI\\\midrule
Disease-pooled vs specialist & \result{P1_auc}\\
Outcome-pooled vs specialist & \result{P2_auc}\\
Joint vs specialist & \result{P5_auc}\\
Joint vs outcome-pooled & \result{P3_auc}\\
Joint vs disease-pooled & \result{P4_auc}\\
Pooling interaction & \result{P6_auc}\\
OPERA (local) vs joint & \result{R1_auc}\\
Original router vs revised & \result{R_orig_auc}\\
OPERA (local) vs per-task XGBoost & \result{O1_auc}\\
Joint vs pooled XGBoost & \result{T6_auc}\\
Pooled vs per-task XGBoost & \result{T7_auc}\\
XGBoost vs logistic regression & \result{T4_auc}\\\bottomrule
\end{tabular}\end{center}

Sharing across diseases produced a larger AUC gain than sharing across
outcomes. Joint learning did not improve discrimination over disease pooling
alone. Routed corrections improved the joint model modestly. The current
local OPERA model remained below XGBoost in both discrimination and
probability accuracy. Its IPA difference from XGBoost was
\result{O1_ipa}, compared with \result{R1_ipa} against the joint model.
Outcome pooling's IPA difference from specialists was \result{P2_ipa}.

At one year, OPERA had higher AUC than XGBoost in 189 of 600 cells, and
higher AUC than joint in 351 of 600. Joint exceeded specialists in 453 of
600. These are counts of higher point estimates, not counts of established
benefit. With horizons pooled within cells, the corresponding counts were
212 of 675, 390 of 675 and 523 of 675.

\begin{center}\small
\begin{tabular}{lr}\toprule
Model & Mean cell AUC, one year\\\midrule
Specialist sequence model & 0.5942\\
Disease-pooled sequence model & 0.7106\\
Outcome-pooled sequence model & 0.6291\\
Joint sequence model & 0.7075\\
OPERA (local, revised router) & 0.7183\\
Original router & 0.7180\\
Per-task XGBoost & 0.7498\\
Pooled XGBoost & 0.7543\\
Logistic regression & 0.6675\\\bottomrule
\end{tabular}\end{center}
These descriptive means use the same 600 cells. They are not a pooled-patient
AUC, and no additional confidence intervals have been calculated for them.

\section*{Routing and support}
Historical diagnostics for the same seed-42 fits show effective-source
summaries of 1.6 for the original router and 70.8 for the revised router.
Overall survival was the original router's top source in 99.8\% of rows.
The revised router's most frequent top source was DLBCL, in 14.7\%.
These diagnose allocation, not biological selectivity or replication across
seeds. A second-seed health report is still required for a stability claim.

The original router's small all-horizon AUC advantage was
\result{R_orig_auc}. Its IPA difference from the revised router was
\result{R_orig_ipa}. At one year, the revised model had a slightly higher
descriptive mean AUC. The comparison depends on horizon and weighting.

Equal-cell own-only versus global-only AUC was
\result{router_own_AB_auc}. Full versus own-only was
\result{router_borrow_AB_auc}, and full versus global-only was
\result{router_both_AB_auc}. Disease-balanced estimates were respectively
\result{router_own_AB_auc_disease_balanced},
\result{router_borrow_AB_auc_disease_balanced} and
\result{router_both_AB_auc_disease_balanced}. These are separately trained
controls. Fixed-fit source removal measures another quantity.

The mean-ablation estimate of dependence on other-disease sources was
0.00459 AUC (95\% interval $-$0.00179 to 0.01039), retaining outcome,
global and shared components. Single-source effects were strongly shrunken
towards zero. The saved shrinkage and unit conversions were reproduced.
There is no matched retrained uniform-mixture control in the supplied bundle.

\begin{center}\small
\begin{tabular}{rrrrr}\toprule
Minimum cases/controls & Entries & Disease pooling & OPERA vs joint & OPERA vs XGBoost\\\midrule
5 & 3,004 & 0.1197 & 0.0087 & $-$0.0280\\
10 & 1,802 & 0.0866 & 0.0157 & $-$0.0356\\
20 & 1,131 & 0.0708 & 0.0149 & $-$0.0348\\
50 & 451 & 0.0233 & 0.0089 & $-$0.0386\\\bottomrule
\end{tabular}\end{center}
These sensitivity means use changing task populations and have no newly
computed intervals. They do not estimate the effect of adding events to the
same task, and do not replace the primary support rule.

Disease-course routing gain was $-$0.01797 across 261 cell--horizon entries.
Excluding 12 Richter-transformation entries diagnostically gave $-$0.00026.
The excluded entries had 5--7 cases and 6--8 controls and repeated three
nested outcomes over four early horizons. They remain in the primary result.
The original router's disease-course mean was also negative. Cardiovascular
routing gain was $-$0.01990 across 46 entries, with only one supported cell at
30 days and two at 60 days. Support and task composition need to accompany
these domain plots.

\section*{Results still required}
\begin{itemize}
\item Refreshed clinical indices with native categories, explicit score/model
identity, training-derived cutoffs, achieved evaluation event rates and
specificity, counts, confusion metrics and paired uncertainty. The old
clinical figures cannot establish the revised clinical headline.
\item National joint and OPERA variants, adaptation and label-efficiency arms.
\item The frozen tuning-selected strategy evaluated on an independent sample.
\item Fresh attribution summary and concept coverage, and the
7/14/21/28-day record-gap analysis.
\item Registry-variable comparator results, any outstanding weighting
sensitivity, and a fit inventory for computational-cost and portfolio claims.
\end{itemize}
The full manuscript retains placeholders for these findings. A missing result
is not a null result. In particular, superiority over every clinical index
and benefit from national pretraining remain unestablished here.

\auditfigure{fig_headline}{Headline comparison: current tuning inputs}
\auditfigure{fig_pooling}{Controlled pooling comparisons}
\auditfigure{supp_horizons}{Clinical domains and horizons}
\auditfigure{fig_hierarchy}{Corrections and training support}
\auditfigure{fig_borrowing_map_mean}{Source removal: mean replacement}
\auditfigure{fig_borrowing_map_zero}{Source removal: zero replacement}
\end{document}
